\documentclass[varwidth=14cm,border=10pt]{standalone}
\usepackage{amsmath,xcolor,array}
\usepackage{fontspec,unicode-math}
\setmainfont{DejaVu Sans}[ItalicFont={DejaVu Sans},ItalicFeatures={FakeSlant=0.2},BoldItalicFont={DejaVu Sans Bold},BoldItalicFeatures={FakeSlant=0.2}]
\setmathfont{texgyredejavu-math.otf}
\setmathfont{DejaVu Sans}[range={up,bfup}]
\setmathfont{DejaVu Sans}[range={it,bfit},FakeSlant=0.2]
\setmathfont{texgyredejavu-math.otf}[range={"2202,"1D715,"03B5,"03F5,"2112,"2131}]% keep partial, epsilon, script caps
\pagecolor[HTML]{1E1E2E}\color[HTML]{F4F4F4}
\definecolor{accent}{HTML}{89B4FA}
\definecolor{warn}{HTML}{F38BA8}
\definecolor{good}{HTML}{A6E3A1}
\newcommand{\heading}[1]{\par\medskip{\large\bfseries\color{accent}#1}\par\smallskip}
\setlength{\parindent}{0pt}
\raggedright
\begin{document}

{\Large\bfseries L08 — Stochastic optimisation and reliability-based design}

\heading{0. The idea in one line}
A deterministic optimum sits {\itshape exactly} on its binding constraints. If the data in
those constraints is uncertain, the optimum holds about half the time. $\beta=0$ is a coin flip.

\heading{1. The five steps}
\begin{enumerate}
\item Write every uncertain constraint as $g_j(x,y)\le0$; list each random $y_i$ with
      $\bar y_i$ and $\sigma_{y_i}$.
\item Write $\bar g_j(x)$ and $\sigma_{g_j}(x)$.
\item \textit{How reliable is a given $x$?} \quad
      $\beta_j = -\bar g_j/\sigma_{g_j}$, \quad $P[g_j\le0]=\Phi(\beta_j)$.
\item \textit{Design for reliability $p_j$:} \quad
      $\bar g_j(x) + \Phi^{-1}(p_j)\,\sigma_{g_j}(x)\le 0$ \quad — fmincon.
\item Check by Monte Carlo: each constraint, {\itshape and the design as a whole}.
\end{enumerate}

\heading{2. Step 2 — the propagation rule}
For a weighted sum of independent random variables, $g=\sum_k \alpha_k y_k$:
\[
  \boxed{\;\bar g = \sum_k \alpha_k \bar y_k\;}
  \qquad
  \boxed{\;\sigma_g^{2} = \sum_k \alpha_k^{2}\,\sigma_{y_k}^{2}\;}
\]
{\bfseries The mean takes the weights — the variance takes the weights {\itshape squared}.}
Signs disappear in the variance: a $-1$ weight contributes $+\sigma^2$.

Worked example, $g = a_2/x - a_1$ (load/size minus strength), weights $1/x$ and $-1$:
\[
  \bar g(x) = \frac{\bar a_2}{x}-\bar a_1,
  \qquad
  \sigma_g(x) = \sqrt{\sigma_{a_1}^{2}+\left(\frac{\sigma_{a_2}}{x}\right)^{2}} .
\]
Both are functions of $x$ alone — the histogram without ever sampling.

\heading{3. Step 3 — the reliability index}
\[
  P[g\le0] = \Phi(\beta),\qquad \beta = -\frac{\bar g}{\sigma_g}
\]
$\beta$ is ``how many standard deviations of safety margin you hold''. It is dimensionless, so
it compares across constraints with different units.
\begin{center}
\begin{tabular}{lccccc}
$\beta$ & 0 & 1 & 2.33 & 3 & \\
$\Phi(\beta)$ & 50\,\% & 84\,\% & 99\,\% & 99.87\,\% & \texttt{normcdf(beta)}\\
\end{tabular}
\end{center}

\heading{4. Step 4 — a probability becomes a constraint}
This is the key algebraic move of the lecture:
\[
  P[g\le0]\ge p
  \;\Longleftrightarrow\;
  \Phi(\beta)\ge p
  \;\overset{\Phi\ \text{increasing}}{\Longleftrightarrow}\;
  -\frac{\bar g(x)}{\sigma_g(x)} \ge \Phi^{-1}(p)
  \;\Longleftrightarrow\;
  \boxed{\;\bar g(x)+\Phi^{-1}(p)\,\sigma_g(x)\le 0\;}
\]
The last step multiplies by $\sigma_g>0$, so the inequality direction survives. The random
variables are {\itshape gone}: only known means and standard deviations remain, and $x$ is the
only unknown. It is now an ordinary L04 nonlinear constraint, \texttt{c <= 0}.
\begin{center}
\begin{tabular}{lccccc}
$p$ & 50\,\% & 90\,\% & 95\,\% & 99\,\% & 99.9\,\%\\
$\Phi^{-1}(p)$ \;\texttt{norminv(p)} & 0.000 & 1.282 & 1.645 & 2.326 & 3.090\\
\end{tabular}
\end{center}
At $p=0.5$, $\Phi^{-1}=0$ and the constraint collapses to $\bar g\le0$ — {\bfseries the
deterministic problem is the $p=50\,\%$ special case}. That is the assertion to put in your
script.

\heading{5. The price of reliability}
\[
  \text{size } x:\quad 0.400\;(50\%)\;\to\;0.506\;(95\%)\;\to\;0.560\;(99\%)\;\to\;0.631\;(99.9\%)
\]
A factor $1.40$ at 99\,\% — which is the safety factor $1.4$–$1.5$ written into civil and
structural design codes, derived rather than decreed. The cost is nonlinear: the last few
nines are the expensive ones, because $\Phi^{-1}$ steepens without bound.

\heading{6. Step 5 — Monte Carlo, and the limitation that matters}
\[
  \texttt{p = mean(g(x,Y) <= 0)}\qquad Y \sim \text{sampled } N \text{ times}
\]
\textcolor{warn}{Each constraint at $p$ is {\itshape not} the design at $p$.} The design needs
{\bfseries all} constraints at once, so the system reliability is lower than any individual
one. With $m$ constraints:
\[
  P[\text{all hold}] \;\le\; \min_j p_j,
  \qquad\text{and if the } g_j \text{ are independent},\quad
  P[\text{all hold}] = \prod_{j=1}^{m} p_j .
\]
Nine constraints each at $99\,\%$, independent: $0.99^{9}=91.4\,\%$ — not $99\,\%$. To target a
{\itshape system} reliability $p_{\text{sys}}$ directly, constrain the product (in logs, for
conditioning):
\[
  \sum_{j=1}^{m}\ln\Phi\!\left(\beta_j(x)\right) \;\ge\; \ln p_{\text{sys}}
  \qquad\Longleftrightarrow\qquad
  \texttt{c} = \ln p_{\text{sys}} - \sum_j \ln\Phi(\beta_j(x)) \le 0
\]
This is one nonlinear constraint coupling all $m$ events, and it is {\itshape exact} whenever
each $g_j$ contains its own distinct random variables.

\heading{7. Other limitations to state honestly}
\begin{enumerate}
\item {\bfseries Normality.} $\Phi$ assumes $g$ is Gaussian. It is exactly so for a linear
      combination of Gaussians; for a nonlinear $g$ (like $a_2/x$) it is a first-order
      approximation — good near the mean, worse far out in the tail, which is precisely where
      high reliability lives. Monte Carlo is the check.
\item {\bfseries Independence.} Correlated $y_i$ add covariance terms to $\sigma_g^2$.
\item {\bfseries The $\sigma$'s are assumptions too.} A CoV of $0.10$ is a modelling choice;
      report it as one.
\end{enumerate}

\end{document}
